% paper-skeleton.tex -- dfact-lcm-square paper lane (2026-10-07). This file
% is the SOURCE of paper.tex: the derived paper.tex is produced from it by
% audit/build.py, which substitutes the @@B:<name>@@ markers with byte-exact
% contiguous excerpts of the shipped proof files and the @@N:<key>@@ markers
% with figures measured at build time. The fact-closure gate scans this
% skeleton (not the derived file), so every fact-type figure here is a marker
% or a spelled-out word and no hand-copied number can go stale. Line ranges of
% the excerpts are recorded in audit/listing-verify.txt, not restated here.
% The assembly sources are the comment-cleaned copies under zenodo/proofs/
% (their code regions are byte-identical to the frozen development of record;
% correspondence record: zenodo/audit/comment-elision-record.md).
% Companion claim of record: papers/dfact-lcm-square's sibling deposit
% claims/dfact-lcm-square (Zenodo DOI reserved at upload).
\documentclass[11pt]{article}

\usepackage[T1]{fontenc}
\usepackage[utf8]{inputenc}
\usepackage{amsmath,amssymb,amsthm}
\usepackage{graphicx}
\usepackage{hyperref}
\usepackage{xurl}
\setlength{\emergencystretch}{3em}
% lstlean.tex — Lean style for the listings package.
% Lineage: Jeremy Avigad (2015), adapted from Assia Mahboubi's SSR style;
% inspired by Olivier Verdier's unixode.sty for the Unicode literate table.
% No CTAN package or upstream repo exists; papers carry this file in their
% source package (cap set arXiv:1907.01449, sphere eversion arXiv:2210.07746).
% Adaptation note: sorry/admit/sorryAx are highlighted in a dedicated error
% tier so that incomplete proofs stand out visually in listings.
\usepackage{listings}
\usepackage{xcolor}

\definecolor{keywordcolor}{rgb}{0.7, 0.1, 0.1}   % red keywords
\definecolor{tacticcolor}{rgb}{0.0, 0.1, 0.6}    % blue tactics
\definecolor{commentcolor}{rgb}{0.4, 0.4, 0.4}   % grey comments
\definecolor{symbolcolor}{rgb}{0.0, 0.1, 0.6}    % blue symbols
\definecolor{sortcolor}{rgb}{0.1, 0.5, 0.1}      % green sorts
\definecolor{errorcolor}{rgb}{0.6, 0, 0}         % dark red: sorry/admit
\definecolor{stringcolor}{rgb}{0.5, 0.1, 0.5}    % purple strings

\lstdefinelanguage{lean}{
  % Anything between $ becomes LaTeX math mode
  mathescape=false,
  % Comments may or not include Latex commands
  texcl=false,
  % keywords, list taken from lean-syntax.txt
  morekeywords=[1]{import, prelude, protected, private, noncomputable,
    definition, meta, renaming, hiding, parameter, parameters, begin,
    constant, constants, lemma, variable, variables, theory, print,
    theorem, example, open, section, namespace, universe, universes, end,
    modifier, modifiers, using, export, exposing, axiom, inductive,
    coinductive, with, structure, class, instance, attribute, local, where,
    by, have, show, let, in, if, then, else, match, do, for, fun, assume,
    from, take, calc, include, omit, infix, infixl, infixr, notation,
    macro, syntax, elab, set_option, initialize, deriving, opaque, abbrev,
    mutual, partial, scoped, termination_by, decreasing_by},
  % Sorts
  morekeywords=[2]{Sort, Type, Prop},
  % tactics, list taken from lean-syntax.txt (tactic keywords)
  morekeywords=[3]{intro, intros, apply, exact, refine, rfl, simp, simp_all,
    simpa, simp_rw, rw, rewrite, erewrite, ring, ring_nf, omega, decide,
    norm_num, norm_cast, induction, cases, rcases, obtain, constructor,
    ext, funext, conv, congr, field_simp, linarith, nlinarith, positivity,
    tauto, trivial, aesop, use, by_cases, by_contra, push_neg, contrapose,
    symm, trans, assumption, contradiction, exfalso, specialize,
    generalize, revert, clear, rename_i, subst, unfold, delta, change,
    convert, split_ifs, interval_cases, fin_cases, zify, gcongr,
    nth_rewrite, suffices, generalizing, at, only, native_decide},
  % warnings - these three tokens mark an unfinished proof; they must never appear at a code position in a delivered file
  morekeywords=[4]{sorry, admit, sorryAx},
  literate=
    {α}{{\ensuremath{\alpha}}}1
    {β}{{\ensuremath{\beta}}}1
    {γ}{{\ensuremath{\gamma}}}1
    {δ}{{\ensuremath{\delta}}}1
    {ε}{{\ensuremath{\varepsilon}}}1
    {ζ}{{\ensuremath{\zeta}}}1
    {η}{{\ensuremath{\eta}}}1
    {θ}{{\ensuremath{\theta}}}1
    {ι}{{\ensuremath{\iota}}}1
    {κ}{{\ensuremath{\kappa}}}1
    {λ}{{\ensuremath{\lambda}}}1
    {μ}{{\ensuremath{\mu}}}1
    {ν}{{\ensuremath{\nu}}}1
    {ξ}{{\ensuremath{\xi}}}1
    {π}{{\ensuremath{\pi}}}1
    {ρ}{{\ensuremath{\rho}}}1
    {σ}{{\ensuremath{\sigma}}}1
    {τ}{{\ensuremath{\tau}}}1
    {υ}{{\ensuremath{\upsilon}}}1
    {φ}{{\ensuremath{\varphi}}}1
    {χ}{{\ensuremath{\chi}}}1
    {ψ}{{\ensuremath{\psi}}}1
    {ω}{{\ensuremath{\omega}}}1
    {Γ}{{\ensuremath{\Gamma}}}1
    {Δ}{{\ensuremath{\Delta}}}1
    {Θ}{{\ensuremath{\Theta}}}1
    {Λ}{{\ensuremath{\Lambda}}}1
    {Ξ}{{\ensuremath{\Xi}}}1
    {Π}{{\ensuremath{\Pi}}}1
    {Σ}{{\ensuremath{\Sigma}}}1
    {Φ}{{\ensuremath{\Phi}}}1
    {Ψ}{{\ensuremath{\Psi}}}1
    {Ω}{{\ensuremath{\Omega}}}1
    {ℵ}{{\ensuremath{\aleph}}}1
    {∀}{{\ensuremath{\forall}}}1
    {∃!}{{\ensuremath{\exists}!}}1
    {∃}{{\ensuremath{\exists}}}1
    {∈}{{\ensuremath{\in}}}1
    {∉}{{\ensuremath{\notin}}}1
    {∘}{{\ensuremath{\circ}}}1
    {∩}{{\ensuremath{\cap}}}1
    {∪}{{\ensuremath{\cup}}}1
    {⊆}{{\ensuremath{\subseteq}}}1
    {⊂}{{\ensuremath{\subset}}}1
    {⊇}{{\ensuremath{\supseteq}}}1
    {⊃}{{\ensuremath{\supset}}}1
    {⊄}{{\ensuremath{\nsubseteq}}}1
    {∅}{{\ensuremath{\emptyset}}}1
    {∧}{{\ensuremath{\wedge}}}1
    {∨}{{\ensuremath{\vee}}}1
    {¬}{{\ensuremath{\neg}}}1
    {⊤}{{\ensuremath{\top}}}1
    {⊥}{{\ensuremath{\bot}}}1
    {↔}{{\ensuremath{\leftrightarrow}}}1
    {→}{{\ensuremath{\rightarrow}}}1
    {←}{{\ensuremath{\leftarrow}}}1
    {↑}{{\ensuremath{\uparrow}}}1
    {⇒}{{\ensuremath{\Rightarrow}}}1
    {⇐}{{\ensuremath{\Leftarrow}}}1
    {⟹}{{\ensuremath{\Longrightarrow}}}1
    {↦}{{\ensuremath{\mapsto}}}1
    {≤}{{\ensuremath{\leq}}}1
    {≥}{{\ensuremath{\geq}}}1
    {≠}{{\ensuremath{\neq}}}1
    {≈}{{\ensuremath{\approx}}}1
    {≡}{{\ensuremath{\equiv}}}1
    {≃}{{\ensuremath{\simeq}}}1
    {≅}{{\ensuremath{\cong}}}1
    {×}{{\ensuremath{\times}}}1
    {·}{{\ensuremath{\cdot}}}1
    {±}{{\ensuremath{\pm}}}1
    {⊕}{{\ensuremath{\oplus}}}1
    {⊗}{{\ensuremath{\otimes}}}1
    {⋆}{{\ensuremath{\star}}}1
    {▸}{{\ensuremath{\blacktriangleright}}}1
    {⟨}{{\ensuremath{\langle}}}1
    {⟩}{{\ensuremath{\rangle}}}1
    {⦃}{{\ensuremath{\{\!\mid}}}1
    {⦄}{{\ensuremath{\mid\!\}}}}1
    {⟦}{{\ensuremath{[\![}}}1
    {⟧}{{\ensuremath{]\!]}}}1
    {⌊}{{\ensuremath{\lfloor}}}1
    {⌋}{{\ensuremath{\rfloor}}}1
    {⌈}{{\ensuremath{\lceil}}}1
    {⌉}{{\ensuremath{\rceil}}}1
    {√}{{\ensuremath{\surd}}}1
    {∣}{{\ensuremath{\mid}}}1
    {∤}{{\ensuremath{\nmid}}}1
    {∎}{{\ensuremath{\blacksquare}}}1
    {⋯}{{\ensuremath{\cdots}}}1
    {…}{{\ensuremath{\ldots}}}1
    {∑}{{\ensuremath{\sum}}}1
    {∏}{{\ensuremath{\prod}}}1
    {⁻¹}{{\ensuremath{^{-1}}}}1
    {¹}{{\ensuremath{^1}}}1
    {²}{{\ensuremath{^2}}}1
    {³}{{\ensuremath{^3}}}1
    {ⁿ}{{\ensuremath{^n}}}1
    {ᵐ}{{\ensuremath{^m}}}1
    {₀}{{\ensuremath{_0}}}1
    {₁}{{\ensuremath{_1}}}1
    {₂}{{\ensuremath{_2}}}1
    {₃}{{\ensuremath{_3}}}1
    {₄}{{\ensuremath{_4}}}1
    {₅}{{\ensuremath{_5}}}1
    {ₙ}{{\ensuremath{_n}}}1
    {ₘ}{{\ensuremath{_m}}}1
    {ₖ}{{\ensuremath{_k}}}1
    {ᵢ}{{\ensuremath{_i}}}1
    {ⱼ}{{\ensuremath{_j}}}1
    {ℤ}{{\ensuremath{\mathbb{Z}}}}1
    {ℕ}{{\ensuremath{\mathbb{N}}}}1
    {ℚ}{{\ensuremath{\mathbb{Q}}}}1
    {ℝ}{{\ensuremath{\mathbb{R}}}}1
    {ℂ}{{\ensuremath{\mathbb{C}}}}1
    {𝔽}{{\ensuremath{\mathbb{F}}}}1
    {𝒫}{{\ensuremath{\mathcal{P}}}}1
    {∗}{{\ensuremath{\ast}}}1
    {•}{{\ensuremath{\bullet}}}1
    {‼}{{!!}}1,
  % Comments
  morecomment=[l]{--},
  morecomment=[s]{/-}{-/},
  morestring=[b]",
  stringstyle=\color{stringcolor},
  keepspaces=true,
  keywordstyle=[1]\color{keywordcolor}\bfseries,
  keywordstyle=[2]\color{sortcolor}\bfseries,
  keywordstyle=[3]\color{tacticcolor}\bfseries,
  keywordstyle=[4]\color{errorcolor}\bfseries\underline,
  escapeinside={(*@}{@*)},
  columns=fullflexible,
  basicstyle=\ttfamily\small,
  commentstyle=\color{commentcolor}\itshape,
  breaklines=true,
  showstringspaces=false,
  frame=single,
  framesep=2mm,
  xleftmargin=4mm,
  numbers=left,
  numberstyle=\tiny\color{commentcolor},
  numbersep=6pt,
}

% Inline Lean code in prose: \lean{Int.ModEq.pow}
\newcommand{\lean}[1]{\lstinline[language=lean,basicstyle=\ttfamily\normalsize]|#1|}

\lstset{language=lean}

% --- XeTeX rendering patch (2026-10-07) --- tectonic/XeTeX does not apply
% the listings literate table to multi-byte characters, so every non-ASCII,
% non-CJK codepoint that occurs in THIS paper's listings is mapped to a math
% glyph as an active character via the standard \lccode/\lowercase idiom
% (ASCII-only source, no extra packages). The mapped set is verified against
% the spliced blocks by audit/build.py at assembly time -- an unmapped
% codepoint aborts the build. Deliberately NOT mapped: U+00B7, U+2014, U+2026
% and all CJK ideographs / CJK punctuation --- the three writing characters
% render through the listings literate table of lstlean.tex under both engines,
% and the CJK ranges route through xeCJK + SimSun (making them active first
% aborts the run, observed in earlier packages). Rendering only: listing
% source bytes are untouched, so the byte-exactness checks are unaffected, and
% the \ifdefined guard leaves the pdflatex path identical to the template.
\ifdefined\XeTeXversion
  {\lccode`\~="00AC \lowercase{\global\catcode`~=\active \global\def~{\ensuremath{\neg}}}}
  {\lccode`\~="203C \lowercase{\global\catcode`~=\active \global\def~{!!}}}
  {\lccode`\~="2115 \lowercase{\global\catcode`~=\active \global\def~{\ensuremath{\mathbb{N}}}}}
  {\lccode`\~="2192 \lowercase{\global\catcode`~=\active \global\def~{\ensuremath{\rightarrow}}}}
  {\lccode`\~="2194 \lowercase{\global\catcode`~=\active \global\def~{\ensuremath{\leftrightarrow}}}}
  {\lccode`\~="2200 \lowercase{\global\catcode`~=\active \global\def~{\ensuremath{\forall}}}}
  {\lccode`\~="2203 \lowercase{\global\catcode`~=\active \global\def~{\ensuremath{\exists}}}}
  {\lccode`\~="2208 \lowercase{\global\catcode`~=\active \global\def~{\ensuremath{\in}}}}
  {\lccode`\~="2223 \lowercase{\global\catcode`~=\active \global\def~{\ensuremath{\mid}}}}
  {\lccode`\~="2227 \lowercase{\global\catcode`~=\active \global\def~{\ensuremath{\wedge}}}}
  {\lccode`\~="2228 \lowercase{\global\catcode`~=\active \global\def~{\ensuremath{\vee}}}}
  {\lccode`\~="2260 \lowercase{\global\catcode`~=\active \global\def~{\ensuremath{\neq}}}}
  {\lccode`\~="2264 \lowercase{\global\catcode`~=\active \global\def~{\ensuremath{\leq}}}}
  {\lccode`\~="27E8 \lowercase{\global\catcode`~=\active \global\def~{\ensuremath{\langle}}}}
  {\lccode`\~="27E9 \lowercase{\global\catcode`~=\active \global\def~{\ensuremath{\rangle}}}}
\fi
\ifdefined\XeTeXversion
  \usepackage{xeCJK}
  \setCJKmainfont{SimSun}
\fi

\newtheorem{theorem}{Theorem}
\newtheorem{lemma}[theorem]{Lemma}
\newtheorem{corollary}[theorem]{Corollary}
\theoremstyle{definition}
\newtheorem{definition}[theorem]{Definition}

\title{The Diophantine equation $x^2 = n!! \cdot \operatorname{lcm}(1,\dots,n)$: the complete solution set,\\ formalized and machine-checked in Lean 4}
\author{Aurora0134\thanks{Contact: via the GitHub account Aurora0134.}}
\date{October 7, 2026}

\begin{document}
\maketitle

\begin{abstract}
We determine the complete solution set of the Diophantine equation
$x^2 = n!! \cdot \operatorname{lcm}(1,\dots,n)$, where $n!!$ is the double
factorial and $\operatorname{lcm}(1,\dots,n)$ the least common multiple of the
integers from $1$ to $n$: for $n \in \mathbb{N}$ the product is a perfect
square if and only if $n \in \{0,1,2,5,7\}$, with square roots
$x = 1, 1, 2, 30, 210$. This settles the second of the two cases that
Novakovi\'c (2026) recorded as believed-finite with the words ``at the moment
we were not able to give an argument''; the first case, $x^2 = n! \cdot n\#$,
was settled separately \cite{nfactpaper}. The answer given here is an
if-and-only-if characterization and is therefore strictly stronger than the
finiteness being asked about. The proof combines a valuation criterion
($n!! \cdot \operatorname{lcm}(1,\dots,n)$ is a square if and only if
$\nu_p(n!!) + \lfloor \log_p n \rfloor$ is even for every prime $p$), an
exclusion of even $n$ by Bertrand's postulate, an exclusion of odd
$n \ge 33$ by an interval-prime lemma for $(n/4, n/3]$ proved in this work as
a prime-certificate chain plus a factorial-ratio tail estimate, and explicit
non-square witnesses for the remaining band. Every statement is
machine-checked in Lean~4 against mathlib v4.34.0: the statements were frozen
by hash before proving, the proof compiles with no \texttt{sorry}, and
\texttt{\#print axioms} reports only \{propext, \texttt{Quot.sound},
\texttt{Classical.choice}\}. The results were developed with AI assistance
under human direction; the Lean 4 kernel is the sole adjudicator, and the
disclosure section states the terms. Lean
sources, frozen statements, and audit logs accompany the paper as a deposited
artifact, and a companion claim of record is deposited separately.
\end{abstract}

\section{Introduction}\label{sec:intro}

The equation of the title asks when the product of the double factorial and
the least common multiple of the initial segment is a perfect square. It
appears, together with the companion equation $x^2 = n! \cdot n\#$, in
Novakovi\'c's 2026 study \cite{novakovic2026} of the Diophantine equation
$P(x) = \prod_{i=1}^{r} H_{n_i}$, where the author writes (verbatim wording;
mathematical notation retypeset): ``We believe that finitely many solutions
occur for $x^2 = n! \cdot n\#$ or $x^2 = n!! \cdot [1, \ldots, n]$, but at the
moment we were not able to give an argument''. In that paper's notation
$[1, \ldots, n]$ denotes $\prod_{p \le n} p^{\lfloor \log_p n \rfloor}$,
which is exactly $\operatorname{lcm}(1,\dots,n)$. The present paper answers
the second of the two cases completely: the equation
$x^2 = n!! \cdot \operatorname{lcm}(1,\dots,n)$ has, for
$n \in \mathbb{N}$, exactly the solutions $n \in \{0,1,2,5,7\}$. The first
case is not claimed here; it was settled separately \cite{nfactpaper}.

Three contributions are recorded together. First, the mathematical result
itself, as an if-and-only-if characterization with explicit witnesses
(Section~\ref{sec:results}); an exclusivity review found no occupying record
for this statement on the channels it could reach
(Section~\ref{sec:related}). Second, a complete machine-checked formalization:
three statements frozen by hash before proving, proved in Lean~4 against
mathlib with no placeholder and with axiom dependencies exactly the three
standard ones (Section~\ref{sec:verify}). Third, the artifact is built so that
a reader can re-run the adjudication --- the frozen statement's hash, the
strict compile's exit status and the axiom printout all ship with the deposit
--- rather than take the text's word for it. This is why the paper addresses a
formal-methods readership as well as a number-theoretic one.

\subsection*{Notation}

We write $n!!$ for the double factorial ($n!! = \prod_{1 \le k \le n,\ k \equiv n \ (2)} k$,
so $7!! = 105$ and $0!! = 1$ by the empty-product convention) and
$\operatorname{lcm}(1,\dots,n)$ for the least common multiple of the integers
$1, \ldots, n$ (so $\operatorname{lcm}(1,\dots,0) = 1$); $\nu_p(m)$ denotes
the exponent of the prime $p$ in $m$, and $\lfloor \log_p n \rfloor$ the
largest $k$ with $p^k \le n$. In the formalization the double factorial is
\lean{Nat.doubleFactorial} (notation \lean{n‼}), the least common multiple is
\lean{Nat.lcmUpto}, and $\nu_p(m)$ is \lean{Nat.factorization}. Division in
the interval-prime lemma (Lemma~\ref{lem:B}) is natural-number division, so
the interval reads $\lfloor n/4 \rfloor < p \le \lfloor n/3 \rfloor$; for
prime existence this coincides with the real reading, since $p$ is an integer.

\section{Mathematical background}\label{sec:background}

\begin{definition}\label{def:objects}
For $n \in \mathbb{N}$, the \emph{double factorial} $n!!$ is the product of
the integers $k \le n$ with $k \equiv n \pmod 2$ (so $0!! = 1!! = 1$ and
$5!! = 15$). For a prime $p$ and $m \ge 1$, $\nu_p(m)$ is the exponent of $p$
in the prime factorization of $m$.
\end{definition}

Two classical ingredients enter. The first is Legendre's formula
$\nu_p(n!) = \sum_{i \ge 1} \lfloor n/p^i \rfloor$, which reaches the proof
through mathlib's \lean{Nat.factorization_factorial} together with the two
identities $n! = n!! \cdot (2m)!! = n!! \cdot 2^m \cdot m!$ for odd
$n = 2m+1$ (\lean{Nat.factorial_eq_mul_doubleFactorial},
\lean{Nat.doubleFactorial_two_mul}). The second is Bertrand's postulate
(\lean{Nat.bertrand}). The shape of the argument is: a natural number is a
square if and only if all its prime valuations are even
(\lean{Nat.isSquare_iff_even_factorization}), and for the product at hand the
total exponent at $p$ is $\nu_p(n!!) + \lfloor \log_p n \rfloor$, because
$\operatorname{lcm}(1,\dots,n)$ contributes exactly $\lfloor \log_p n \rfloor$
(\lean{Nat.factorization_lcmUpto}). So it suffices to find, for each $n$
outside $\{0,1,2,5,7\}$, a prime $p$ at which that total is odd. Even $n$ are
killed by Bertrand; odd $n \ge 33$ by an interval prime in $(n/4, n/3]$; the
remaining band is decided pointwise.

\paragraph{Where this sits in the literature.} The equation is the second case
of the remark in \cite{novakovic2026}; the first case $x^2 = n! \cdot n\#$ was
settled separately \cite{nfactpaper}, and its
interval-prime lemma (for $(n/3, n/2]$) is the structural model for
Lemma~\ref{lem:B} below, with the same certificate-chain-plus-tail
architecture. The same author's other papers
\cite{novakovic2025,novakovic2026b,novakovic2026c,novakovic2026d} are recorded
with the restricted wording that their objects do not contain this equation
(each was read in full; see Section~\ref{sec:related}). Adjacent records that
must be distinguished from the present statement --- the squarefree part of
$\operatorname{lcm}(1,\dots,n)$ (OEIS A056839), the value $(p_n\#)^2$ (OEIS
A061742), the squarefree part of $n!$ (OEIS A055204), the disproved
Erd\H{o}s problem \#399 \cite{erdos399}, primorial-core Pell equations
\cite{peixoto2026}, and square products of factorials \cite{yudin2026} --- are
delimited one by one in Section~\ref{sec:related}.

\section{Formalisation setup}\label{sec:setup}

The formalization is written against Lean~4 \cite{lean4}, version
\texttt{v4.34.0} (the \texttt{leanprover/lean4} release of that version),
with mathlib4 \cite{mathlib} \texttt{v4.34.0} at revision \texttt{5ed29652};
the pin ships with the artifact (\texttt{proofs/lean-toolchain}). The three
statements of Listing~\ref{lst:stmt} were frozen by hash before proving
began; the proving stage may replace proof bodies only, and the final audit
re-compares the delivered statement headers against the frozen snapshot under
two whitespace conventions. Two results the final development consumes come
from an accompanying library rather than from mathlib: the interval-prime lemma
of Lemma~\ref{lem:B} and one bridge used by the witness layer. Their module
names appear among the imports of the shipped final file; they are not part of
this deposit (Section~\ref{sec:limits}(b)).

\paragraph{Encoding decisions.} Four encoding choices are worth recording.
(i) The double factorial is mathlib's \lean{Nat.doubleFactorial} and the
least common multiple is \lean{Nat.lcmUpto}; the latter's factorization
identity $\nu_p(\operatorname{lcm}(1,\dots,n)) = \lfloor \log_p n \rfloor$ is
\lean{Nat.factorization_lcmUpto}, which is what makes the valuation criterion
a one-line consequence of the product law. (ii) The main result is delivered
in two forms: the iff form \lean{dfact_lcm_square_solution_set} (``the product
is a square if and only if $n \in \{0,1,2,5,7\}$'') and the
existence-witness form \lean{exists_sq_dfact_lcm_iff} (``$x^2 = n!! \cdot
\operatorname{lcm}(1,\dots,n)$ is solvable if and only if $n \in
\{0,1,2,5,7\}$''), bridged by \lean{isSquare_iff_exists_sq}; the solution set
is written as a disjunction chain, which on $\mathbb{N}$ is equivalent to
membership in a finite set. (iii) The interval-prime lemma is stated with
natural-number division, faithful to the real reading for prime existence as
recorded in the formalization's fidelity note. (iv) The interval-prime lemma
is proved in this work rather than cited: it is a finite certificate chain
covering a bounded segment plus a factorial-ratio estimate for the tail, and
it is not an instance of any published interval theorem
(Section~\ref{sec:limits}(b)).

\paragraph{Audit mechanism.} The statement is frozen by hash; a strict compile
must exit $0$ with no placeholder (\texttt{sorry}/\texttt{admit}) anywhere in
the file, comments included; \texttt{\#print axioms} must report only the
three standard axioms; and the statement comparison is re-run by an
independent audit at the end. Two solution-set perturbation probes (the
statement altered to a superset and to a subset of the true solution set)
must turn the proof red, and they did. The verdicts for this paper are
reproduced in Section~\ref{sec:verify} and ship with the artifact.

\section{Main results}\label{sec:results}

% LEAN: tasks/20261006-dfact-lcm-square :: DFactLcmSquare.dfact_lcm_square_solution_set grade=完全证明
\begin{theorem}[complete solution set, iff form]\label{thm:main}
For every $n \in \mathbb{N}$, $n!! \cdot \operatorname{lcm}(1,\dots,n)$ is a
perfect square if and only if $n = 0$ or $n = 1$ or $n = 2$ or $n = 5$ or
$n = 7$.
\end{theorem}

% LEAN: tasks/20261006-dfact-lcm-square :: DFactLcmSquare.exists_sq_dfact_lcm_iff grade=完全证明
\begin{theorem}[existence-witness form]\label{thm:exists}
For every $n \in \mathbb{N}$, the equation
$x^2 = n!! \cdot \operatorname{lcm}(1,\dots,n)$ has a solution
$x \in \mathbb{N}$ if and only if $n \in \{0,1,2,5,7\}$.
\end{theorem}

% LEAN: tasks/20261006-dfact-lcm-square :: DFactLcmSquare.isSquare_dfact_mul_lcmUpto_iff grade=完全证明
\begin{lemma}[valuation criterion]\label{lem:A}
For every $n \in \mathbb{N}$: $n!! \cdot \operatorname{lcm}(1,\dots,n)$ is a
perfect square if and only if $\nu_p(n!!) + \lfloor \log_p n \rfloor$ is even
for every prime $p$.
\end{lemma}

% LEAN: tasks/20261006-dfact-lcm-square :: DFactLcmSquare.dfact_lcm_witnesses grade=完全证明
\begin{theorem}[solution witnesses]\label{thm:witnesses}
$1^2 = 0!! \cdot \operatorname{lcm}(1,\dots,0)$,
$1^2 = 1!! \cdot \operatorname{lcm}(1,\dots,1)$,
$2^2 = 2!! \cdot \operatorname{lcm}(1,\dots,2)$,
$30^2 = 5!! \cdot \operatorname{lcm}(1,\dots,5)$ and
$210^2 = 7!! \cdot \operatorname{lcm}(1,\dots,7)$.
\end{theorem}

\begin{lstlisting}[caption={Frozen statement snapshot (Lean 4): the complete
frozen set, taken from the statement source shipped as
\texttt{proofs/dfact-lcm-square.statement.lean} (a comment-cleaned edition:
comments carrying pre-publication workflow narrative were blanked with line
numbers preserved; the statement text is byte-identical to the frozen snapshot
of record, sha256 \texttt{9baa672f}; correspondence record
\texttt{audit/comment-elision-record.md}).},label={lst:stmt}]
theorem isSquare_dfact_mul_lcmUpto_iff (n : ℕ) :
    IsSquare (n‼ * Nat.lcmUpto n) ↔
      ∀ p : ℕ, p.Prime → Even ((n‼).factorization p + p.log n) := by

theorem dfact_lcm_witnesses :
    1 ^ 2 = 0‼ * Nat.lcmUpto 0 ∧ 1 ^ 2 = 1‼ * Nat.lcmUpto 1 ∧
      2 ^ 2 = 2‼ * Nat.lcmUpto 2 ∧ 30 ^ 2 = 5‼ * Nat.lcmUpto 5 ∧
        210 ^ 2 = 7‼ * Nat.lcmUpto 7 := by

theorem dfact_lcm_square_solution_set (n : ℕ) :
    IsSquare (n‼ * Nat.lcmUpto n) ↔ n = 0 ∨ n = 1 ∨ n = 2 ∨ n = 5 ∨ n = 7 := by
\end{lstlisting}

\section{Proofs}\label{sec:proofs}

The final proof file (a single Lean module, 333 lines) assembles
five components: the valuation criterion, the witness layer, the two exclusion
arms, and the assembly of the two main theorems. The file ships byte-exactly
in the artifact under \texttt{proofs/}; the excerpts below are byte-identical
contiguous blocks of it (independent mechanical record:
\texttt{audit/listing-verify.txt}, which ships in the deposit). The shipped
file is the comment-cleaned edition; its code is byte-identical to the frozen
development (record: \texttt{audit/comment-elision-record.md}). We describe
the proof in the order in which it is built.

\subsection{The valuation criterion}

The multiplicative spine is the product law for valuations together with the
least-common-multiple factorization identity: for a prime $p$,
\[
\nu_p\bigl(n!! \cdot \operatorname{lcm}(1,\dots,n)\bigr)
  = \nu_p(n!!) + \lfloor \log_p n \rfloor .
\]
A natural number is a square if and only if all its prime valuations are even
(\lean{Nat.isSquare_iff_even_factorization}), so the criterion of
Lemma~\ref{lem:A} follows term by term. Listing~\ref{lst:crit} shows the whole
argument: rewrite the square predicate into the even-valuation form, push the
universal quantifier through the congruence, and reduce the sum to the two
factors.

\begin{lstlisting}[caption={The valuation criterion (Lemma~\ref{lem:A}): the
criterion bridge in full.},label={lst:crit}]
theorem isSquare_dfact_mul_lcmUpto_iff (n : ℕ) :
    IsSquare (n‼ * Nat.lcmUpto n) ↔
      ∀ p : ℕ, p.Prime → Even ((n‼).factorization p + p.log n) := by
  rw [Nat.isSquare_iff_even_factorization]
  refine forall_congr' fun p => imp_congr_right fun hp => ?_
  simp only [Nat.factorization_mul (Nat.doubleFactorial_pos n).ne' (Nat.lcmUpto_ne_zero n),
    Finsupp.add_apply, Nat.factorization_lcmUpto n hp]
\end{lstlisting}

\subsection{The witness layer}

The five solutions are exhibited pointwise
(Theorem~\ref{thm:witnesses}, closed by decision). For the non-solutions, six
explicit factorizations are recorded as separate theorems
\lean{nonsquare\_at\_3}, \lean{nonsquare\_at\_9}, \lean{nonsquare\_at\_11},
\lean{nonsquare\_at\_13}, \lean{nonsquare\_at\_29} and
\lean{nonsquare\_at\_31}: each exhibits a prime $p$ with
$\nu_p(n!! \cdot \operatorname{lcm}(1,\dots,n)) = 2e+1$ odd --- for instance
at $n = 3$ the value is $18 = 2^1 \cdot 9$, and at $n = 9$ it is
$2381400 = 3^5 \cdot 9800$ --- and the library bridge
\lean{not\_isSquare\_of\_odd\_factorization} turns an odd valuation into a
non-square (a square $m = r \cdot r$ would have
$\nu_p(m) = \nu_p(r) + \nu_p(r)$ even). Listing~\ref{lst:wit} shows the first
of the six in full.

\begin{lstlisting}[caption={The first explicit non-square witness (at $n = 3$)
and the witness layer for the five solutions.},label={lst:wit}]
theorem dfact_lcm_witnesses :
    1 ^ 2 = 0‼ * Nat.lcmUpto 0 ∧ 1 ^ 2 = 1‼ * Nat.lcmUpto 1 ∧
      2 ^ 2 = 2‼ * Nat.lcmUpto 2 ∧ 30 ^ 2 = 5‼ * Nat.lcmUpto 5 ∧
        210 ^ 2 = 7‼ * Nat.lcmUpto 7 := by
  decide

theorem nonsquare_at_3 : ¬ IsSquare (3‼ * Nat.lcmUpto 3) := by
  have hM : (3‼ * Nat.lcmUpto 3) = 2 ^ 1 * 9 := by decide
  apply NFactPrimorialSquare.not_isSquare_of_odd_factorization (p := 2) (e := 0)
    (mul_ne_zero (Nat.doubleFactorial_pos 3).ne' (Nat.lcmUpto_ne_zero 3))
  rw [hM, Nat.factorization_mul (by norm_num : ((2:ℕ)^1) ≠ 0) (by norm_num : (9:ℕ) ≠ 0),
    Finsupp.add_apply, Nat.factorization_eq_zero_of_not_dvd (by decide : ¬ (2:ℕ) ∣ 9),
    Nat.factorization_pow, Finsupp.smul_apply,
    Nat.Prime.factorization Nat.prime_two, Finsupp.single_eq_same]; norm_num [nsmul_eq_mul]
\end{lstlisting}

\subsection{The even arm}

For even $n = 2m \ge 4$, Bertrand's postulate supplies a prime
$p \in (m, 2m]$. Then $p$ is odd, $p > m$, and $p \le n < p^2$, so
$\nu_p(n!!) = 0$ (the prime divides neither $2^m$ nor $m!$) and
$\lfloor \log_p n \rfloor = 1$: an odd total. The argument is
Listing~\ref{lst:even} in full; the only library input is Bertrand's
postulate and Legendre's formula in the form
\lean{Nat.factorization\_factorial}.

\begin{lstlisting}[caption={The even arm in full: even $n \ge 4$ are excluded by
Bertrand's postulate.},label={lst:even}]
theorem even_dfact_lcm_nonsquare (n : ℕ) (hn : 4 ≤ n) (h2 : Even n) :
    ¬ IsSquare (n‼ * Nat.lcmUpto n) := by
  obtain ⟨m, rfl⟩ : ∃ m, n = 2 * m := by
    obtain ⟨m, hm⟩ := h2
    exact ⟨m, by omega⟩
  have hm2 : 2 ≤ m := by omega
  obtain ⟨p, hp, hmp, hpn⟩ := Nat.bertrand m (by omega)
  have hpsq : 2 * m < p * p := by
    have h1 : 2 * m < 2 * p := by omega
    have h2' : 2 * p ≤ p * p := Nat.mul_le_mul (by omega : 2 ≤ p) (le_refl p)
    exact lt_of_lt_of_le h1 h2'
  have hndvd2m : ¬ p ∣ 2 ^ m := fun hd => by
    have h1 := hp.dvd_of_dvd_pow hd
    have h2 := Nat.le_of_dvd (by norm_num : (0:ℕ) < 2) h1
    omega
  have hndvdm : ¬ p ∣ m ! := fun hd => by
    have := (Nat.Prime.dvd_factorial hp).mp hd
    omega
  have hfact_dfact : ((2 * m)‼).factorization p = 0 := by
    rw [Nat.doubleFactorial_two_mul,
      Nat.factorization_eq_zero_of_not_dvd (hp.not_dvd_mul hndvd2m hndvdm)]
  have hlog1 : p.log (2 * m) = 1 := by
    have hge : 1 ≤ Nat.log p (2 * m) :=
      Nat.le_log_of_pow_le hp.one_lt (by rw [pow_one]; omega)
    have hlt : Nat.log p (2 * m) < 2 :=
      Nat.log_lt_of_lt_pow (by omega) (by rwa [pow_two])
    omega
  apply NFactPrimorialSquare.not_isSquare_of_odd_factorization (p := p) (e := 0)
    (mul_ne_zero (Nat.doubleFactorial_pos _).ne' (Nat.lcmUpto_ne_zero _))
  rw [Nat.factorization_mul (Nat.doubleFactorial_pos _).ne' (Nat.lcmUpto_ne_zero _),
    Finsupp.add_apply, hfact_dfact, Nat.factorization_lcmUpto _ hp, hlog1]
\end{lstlisting}

\subsection{The odd arm and the interval-prime lemma}

For odd $n$ with a prime $p \in (n/4, n/3]$ and $n < p^2$, the same bookkeeping
gives $\nu_p(n!!) = 2$ and $\lfloor \log_p n \rfloor = 1$, an odd total. The
mechanism is the generic lemma \lean{odd\_dfact\_lcm\_nonsquare\_of\_prime},
whose statement is exactly those hypotheses; the content is
Lemma~\ref{lem:B}, the existence of such a prime for all odd $n \ge 33$.

\begin{lemma}[interval-prime lemma]\label{lem:B}
For every $n \ge 33$ there exists a prime $p$ with $n/4 < p \le n/3$.
\end{lemma}

Lemma~\ref{lem:B} is proved in this work, not cited. Its architecture has two
parts. On a bounded initial segment it is a certificate chain: a single hop
observes that a prime $q$ with $3q \le n \le 4q - 1$ witnesses
$n/4 < q \le n/3$, and a chain of consecutive primes with the continuity
condition $q' \le 4q/3$ covers every $n$ from the start up to a computable
bound. Beyond that bound the argument is a factorial-ratio estimate: if no
prime lies in $(n/4, n/3]$ then the binomial coefficient
$\binom{\lfloor n/3 \rfloor}{\lfloor n/4 \rfloor, \lfloor n/12 \rfloor}$
exceeds what the surrounding interval can supply, a contradiction. The lemma
lives in an accompanying library and is consumed here by import. It is not
reproduced in this paper's listings, because its proof is a separate
development of the same grade; the deposit states the lemma and the architecture
of its proof, and the shipped final file names the imported module. A reader who
wants to re-verify the odd arm therefore needs that library as well, and this
dependency is recorded in the Limitations (Section~\ref{sec:limits}(b)).

\subsection{Assembly}

The main theorem's forward direction is by contradiction, split at $n = 4$:
below the split the three small cases are decided by the witness theorems and
decision procedures; at and above it, the parity of $n$ selects the even or
the odd arm, and the odd arm's remaining band $9 \le n < 33$ is closed by the
six explicit witnesses interleaved with the generic odd-arm lemma at $p = 5$
and $p = 7$. The reverse direction is the witness layer.
Listing~\ref{lst:main} is the opening of that assembly.

\begin{lstlisting}[caption={The main theorem's assembly: the forward direction's
case split, from the contradiction to the two arms.},label={lst:main}]
theorem dfact_lcm_square_solution_set (n : ℕ) :
    IsSquare (n‼ * Nat.lcmUpto n) ↔ n = 0 ∨ n = 1 ∨ n = 2 ∨ n = 5 ∨ n = 7 := by
  constructor
  · intro h
    by_contra hne
    push_neg at hne
    rcases lt_or_ge n 4 with hlt | hge
    · interval_cases n
      · omega
      · omega
      · omega
\end{lstlisting}

\section{Verification}\label{sec:verify}

The final file compiles in strict mode (errors rejected, plus a textual scan
for \texttt{sorry}/\texttt{admit} including comments) with exit status $0$;
the only diagnostic is one style-level deprecation warning on
\texttt{push\_neg}. \texttt{\#print axioms} on each of the twelve theorems
returns exactly \texttt{[propext, Classical.choice, Quot.sound]} ---
reproduced verbatim in Listing~\ref{lst:axioms}, which is the machine output
that ships in the deposit. The statement comparison against the frozen
snapshot is clean declaration by declaration, and the two solution-set
perturbation probes turn the proof red at exactly the perturbed theorem and
its direct downstream. The independent final audit's verdict for this result
is \emph{complete proof} and ships with the artifact.

\begin{lstlisting}[caption={Axiom printout for the three frozen statements and
the accompanying existential form (machine output, verbatim).},label={lst:axioms}]
<source-task>/attempts/axioms-probe.lean:257:4: warning: `push_neg` has been deprecated. Prefer using `push Not` instead.
If you'd rather continue using `push_neg` in your project, you can implement it as follows:
```
open Lean.Parser.Tactic in
macro "push_neg" cfg:optConfig loc:(location)? : tactic =>
  `(tactic| push $cfg:optConfig Not $[$loc]?)
```
'DFactLcmSquare.isSquare_dfact_mul_lcmUpto_iff' depends on axioms: [propext, Classical.choice, Quot.sound]
'DFactLcmSquare.dfact_lcm_witnesses' depends on axioms: [propext, Classical.choice, Quot.sound]
'DFactLcmSquare.nonsquare_at_3' depends on axioms: [propext, Classical.choice, Quot.sound]
'DFactLcmSquare.nonsquare_at_9' depends on axioms: [propext, Classical.choice, Quot.sound]
'DFactLcmSquare.nonsquare_at_11' depends on axioms: [propext, Classical.choice, Quot.sound]
'DFactLcmSquare.nonsquare_at_13' depends on axioms: [propext, Classical.choice, Quot.sound]
'DFactLcmSquare.nonsquare_at_29' depends on axioms: [propext, Classical.choice, Quot.sound]
'DFactLcmSquare.nonsquare_at_31' depends on axioms: [propext, Classical.choice, Quot.sound]
'DFactLcmSquare.odd_dfact_lcm_nonsquare_of_prime' depends on axioms: [propext, Classical.choice, Quot.sound]
'DFactLcmSquare.even_dfact_lcm_nonsquare' depends on axioms: [propext, Classical.choice, Quot.sound]
'DFactLcmSquare.dfact_lcm_square_solution_set' depends on axioms: [propext, Classical.choice, Quot.sound]
'DFactLcmSquare.exists_sq_dfact_lcm_iff' depends on axioms: [propext, Classical.choice, Quot.sound]
\end{lstlisting}

\section{Related work and novelty}\label{sec:related}

The exclusivity review's four-state adjudication is \emph{no occupying record
found}, with strength labels per channel; the full query record, including
every zero hit, ships as \texttt{audit/novelty-record.md}. The channels and what
they reached: (i)~\emph{OEIS value columns} --- the solution-$n$ column, the
solution-root column $x = 1,1,2,30,210$, and the object value column
$a(n) = n!! \cdot \operatorname{lcm}(1,\dots,n)$ under four prefix variants
(plus alternating-parity sub-columns and drop-prefix cuts in the review pass)
are all null; the residual hits are unrelated items, or A061742 $= (p_n\#)^2$,
which is entered precisely because $x = n\#$ and is not a registration of this
equation. (ii)~\emph{arXiv metadata} --- fourteen load-bearing query forms
with the positive control \texttt{all:"double factorial"} returning
thirty-four hits, confirming the channel was live: no same form at the
metadata layer. (iii)~\emph{Neighbouring full texts} --- seven items read in
full, together with all seventeen references of the anchor: zero coverage of
this equation. (iv)~\emph{Citation axis} --- Semantic Scholar and OpenAlex,
two independent channels, on the anchor: \texttt{cited\_by\_count} $= 0$ on
both, so the anchor's open case has not been closed by anyone citing it.
(v)~\emph{Public Lean ecosystem} --- \texttt{doubleFactorial} occurs in exactly
one file of the whole public-conjectures repository (A001818, a different
object); the statement layer has zero registration. (vi)~\emph{mathlib} ---
ripgrep with the positive control \texttt{Nat.Prime} returning twenty-one
files: \texttt{doubleFactorial} occurs in four files, none in a
number-theoretic square context, \texttt{lcmUpto} in one, and the two objects
never co-occur.

Three delimitations are recorded rather than left implicit. First, the
$\operatorname{lcm}(1,\dots,n)$ side of the object is occupied at the
\emph{sequence} level: OEIS A056839 registers the squarefree part of
$\operatorname{lcm}(1,\dots,n)$; that entry does not mention the double
factorial anywhere, and it registers a sequence, not the characterizing
solution set of this equation. Second, the review layer records that Crossref
is a relevance-ranked channel and that the OpenAlex title search missed the
anchor by indexing status; neither is used as a load-bearing zero hit. Third,
Erd\H{o}s problem \#399 is \emph{not} an open anchor --- it was disproved on
2026-02-04 (counterexample $10! = 48^4 - 36^4$) --- and is mentioned only as a
same-family adjacency; and A055204 asks for the \emph{smallest} $a$ making
$n! \cdot a(n)$ a square, a different object from the present statement.

\section{Limitations and threats to validity}\label{sec:limits}

(a) \textbf{What is proved.} The kernel layer proves exactly the three frozen
statements of Listing~\ref{lst:stmt} plus the proof-side additions recorded in
the audit; the grade of each is \emph{complete proof}. Nothing beyond is
claimed: no growth statement, no other multiplier in place of
$\operatorname{lcm}(1,\dots,n)$, and no first case.

(b) \textbf{The interval-prime lemma is proved here, not cited.}
Lemma~\ref{lem:B} is a certificate chain over a bounded segment plus a
factorial-ratio tail estimate, both built in the accompanying library
development rather than taken from a published interval theorem. The paper
makes no claim that it is an instance of, or comparable in strength to, any
published interval theorem; in particular it is not presented as an
application of Nagura's bound. The corresponding disclosure on the reader's
side: this deposit ships the final development, which imports that lemma
module, but not the module itself, so the odd arm cannot be re-verified from
the deposited files alone. Both the odd-arm mechanism and one bridge used by
the witness layer are library items, each with its own separately checked final
state; their grades are inherited here, not re-established.

(c) \textbf{The first case is not claimed.} The companion equation
$x^2 = n! \cdot n\#$ of the same remark in \cite{novakovic2026} is not claimed
here; it was settled separately \cite{nfactpaper}.

(d) \textbf{No registration claim.} No OEIS registration is claimed or
performed; the review layer found no registration of the value columns, and
any registration action is the author's own.

(e) \textbf{An observation, not a theorem.} The observation that the five
solution roots are $x = n\#$ (the primorials $0\#, 1\#, 2\#, 5\#, 7\# =
1, 1, 2, 30, 210$) is a numerical observation recorded alongside the theorem;
it was not formalized as a statement, and the paper does not present it as a
proved proposition.

(f) \textbf{A numerical coincidence, explicitly disclosed.} The intersection
$\{5,7\}$ of this solution set with the Brocard solutions $\{4,5,7\}$ is a
numerical coincidence; this result implies nothing in either direction about
the Brocard equation, and no Brocard-direction conclusion is claimed.

(g) \textbf{Comment-level deltas.} The shipped \texttt{.lean} files are copies
in which comments carrying pre-publication workflow narrative were blanked
(line numbers preserved, code bytes untouched); per file the comment-free code
view of the source and of the copy are byte-for-byte equal, asserted by the
generating script. Statements and proof terms are therefore exactly the frozen
artifacts, and the compile and axiom records shipped in \texttt{audit/} were
produced from the sources in their original, unblanked form.

(h) \textbf{Semantic fidelity of the formalization.} The theorems are about
the Lean definitions. Their fidelity to the informal reading rests on the
frozen statements being the intended encoding; a semantic round-trip check
(formal statement back to natural language, adjudicated consistent) was
performed at the formalization stage, and a further spot check of this paper's
informal exposition against the frozen statements was done at assembly. A
reader who disagrees with the encoding is disagreeing with the statement, not
with the proof.

(i) \textbf{Novelty review limits.} The zero-hit verdict is a statement about
the channels reached (arXiv metadata, OEIS keyword and value faces, seven
neighbouring full texts, two citation channels, the public Lean ecosystem via a
repository tarball, the local mathlib tree), not about the literature as a
whole; background items recorded as undecided at the full-text level stay in
the source task card, which is not shipped, and none of them is load-bearing
for this result. No ``first formalization'' claim is made.

(j) \textbf{Numeric evidence is layered.} The factorization witnesses quoted
in the text are numeric-level evidence; every theorem-level conclusion is
adjudicated by the kernel and depends on no numeric evidence. The
$[0,64]$ squareness scan reported by the audit is corroboration, not a
premise.

\section*{Data availability}
Lean sources, the frozen statement, compile and axiom records, and the full
novelty-search record are packaged on Zenodo (DOI: \texttt{10.5281/zenodo.23309984}).
The paper is released under CC BY 4.0 and the code under MIT. A companion
claim of record for the same kernel result is deposited separately.

% Disclosure section (2026-10-07 minimal-disclosure ruling): the fixed three
% sentences, where only the model name and the version numbers are placeholders.
% The model name comes from the source task card's node snapshot (the last
% segment of its wire id); no channel prefix, no resolution level and no capture
% time may be copied, and no enumeration of internal stages or mechanism
% self-description words. This section is the paper's substantive AI disclosure;
% the abstract carries one cross-reference sentence, and no other passage
% discusses the development process.
\section*{Use of generative AI}
The formalization and proofs reported in this paper were developed with AI
assistance (large language model: \texttt{kimi-k3}) under human direction and
review. All formal statements and proofs reported in this paper were checked by
the Lean 4 kernel (Lean v4.34.0, mathlib v4.34.0): every proof compiles with
no \texttt{sorry}, and \texttt{\#print axioms} reports only \{propext,
Quot.sound, Classical.choice\}. The author reviewed all content and takes full
responsibility for it.

\section*{Author contributions}
The author selected the problem, verified the semantics of the formal
statements against the informal ones, and reviewed the kernel evidence and the
novelty record before signing off the deposit.

\begin{thebibliography}{99}
\bibitem{novakovic2026}
S.~Novakovi\'c.
\newblock The Diophantine equation $P(x)=\overset{r}{\underset{i=1}{\prod}}H_{n_i}$.
\newblock arXiv:2601.16757, 2026.

\bibitem{novakovic2025}
S.~Novakovi\'c.
\newblock Diophantine equations involving double factorials.
\newblock arXiv:2510.24312, 2025.

\bibitem{novakovic2026b}
S.~Novakovi\'c.
\newblock On a generalization of the Brocard--Ramanujan Diophantine equation.
\newblock arXiv:2602.09687, 2026.

\bibitem{novakovic2026c}
S.~Novakovi\'c.
\newblock Products of factorials which are products of factorials.
\newblock arXiv:2602.23838, 2026.

\bibitem{novakovic2026d}
S.~Novakovi\'c.
\newblock A comment on the equation $n!!=a_1!!\cdots a_t!!$.
\newblock arXiv:2604.09730, 2026.

\bibitem{yudin2026}
F.~Yudin.
\newblock Square products of factorials and a conjecture of Erd\H{o}s and Graham.
\newblock arXiv:2610.01899, 2026.

\bibitem{peixoto2026}
L.~V. Peixoto.
\newblock The Brocard Problem: Structural Invariants, $p$-Adic Density, and the Generative Sieve.
\newblock arXiv:2606.23485, 2026.

\bibitem{erdos399}
Erd\H{o}s problem \#399, erdosproblems.com. Recorded as disproved
(2026-02-04); mentioned here only as a same-family adjacency.

\bibitem{nfactpaper}
A.~Aurora0134.
\newblock The Diophantine equation $x^2 = n! \cdot n\#$: the complete solution
set, formalized and machine-checked in Lean~4.
\newblock Zenodo deposit, 2026. \url{https://doi.org/10.5281/zenodo.23179326}

\bibitem{lean4}
L.~de Moura and S.~Ullrich.
\newblock The Lean 4 theorem prover and programming language.
\newblock \emph{CADE-28}, 2021. \url{https://doi.org/10.1007/978-3-030-79876-5_37}

\bibitem{mathlib}
The mathlib Community.
\newblock The Lean mathematical library.
\newblock \emph{CPP 2020}, 2020. \url{https://doi.org/10.1145/3372885.3373824}
\end{thebibliography}

\end{document}
